\chapter{Tableaux et fonctions en PHP}

\begin{synthese}[Au programme]
Ranger plusieurs valeurs dans un même \textbf{tableau} (indexé ou associatif),
le parcourir avec \code{for} et \code{foreach}, exploiter les fonctions prêtes à
l'emploi (\code{count}, \code{sort}, \code{in\_array}\dots), puis écrire ses
propres \textbf{fonctions} pour organiser et réutiliser son code.
\end{synthese}

\section{Les tableaux indexés}

Jusqu'ici, une variable ne contenait qu'\emph{une seule} valeur. Pour gérer les
notes d'une classe de 40 élèves, il faudrait 40 variables\,! Le \textbf{tableau}
règle ce problème : il stocke plusieurs valeurs sous un seul nom.

\begin{definition}
Un \textbf{tableau} (\emph{array}) est une variable capable de contenir
\textbf{plusieurs valeurs}. Dans un tableau \textbf{indexé}, chaque valeur est
repérée par un numéro de position appelé \textbf{indice}, qui commence à
\textbf{0}.
\end{definition}

\subsection{Créer un tableau}

PHP offre deux écritures équivalentes : la fonction \code{array()} ou les
crochets \code{[]} (plus courte, recommandée).

\begin{codephp}
<?php
  // Écriture classique
  $eleves = array("Awa", "Bello", "Chimene", "Daniel");
  // Écriture courte (identique)
  $villes = ["Yaoundé", "Douala", "Maroua", "Bafoussam"];

  // Un tableau vide, que l'on remplira ensuite
  $notes = [];
?>
\end{codephp}

\begin{syntaxe}
\code{\$tableau = [valeur0, valeur1, valeur2, \dots];}
\end{syntaxe}

\subsection{Accéder à une valeur par son indice}

On écrit le nom du tableau suivi de l'indice entre crochets. Le \textbf{premier}
élément porte l'indice \code{0}, le deuxième \code{1}, etc.

\begin{codephp}
<?php
  $eleves = ["Awa", "Bello", "Chimene", "Daniel"];
  echo $eleves[0];   // affiche : Awa
  echo $eleves[2];   // affiche : Chimene

  // On peut aussi modifier une case
  $eleves[1] = "Bello Ngono";
  echo $eleves[1];   // affiche : Bello Ngono
?>
\end{codephp}

\begin{attention}
Un tableau de 4 éléments a pour indices \code{0}, \code{1}, \code{2} et \code{3}.
Le dernier indice est donc \textbf{4 - 1 = 3}. Accéder à \code{\$eleves[4]}
renvoie une valeur \code{NULL} et déclenche un avertissement.
\end{attention}

\subsection{Ajouter un élément}

Pour ajouter une valeur à la fin d'un tableau, on utilise des crochets
\textbf{vides} \code{[]} : PHP choisit automatiquement le bon indice.

\begin{codephp}
<?php
  $notes = [];          // tableau vide
  $notes[] = 12;        // va à l'indice 0
  $notes[] = 15.5;      // va à l'indice 1
  $notes[] = 9;         // va à l'indice 2
  echo $notes[1];       // affiche : 15.5
?>
\end{codephp}

\begin{remarque}
\code{\$notes[] = 12;} signifie « ajoute 12 à la fin ». On verra plus loin la
fonction \code{array\_push} qui fait la même chose et permet d'ajouter plusieurs
valeurs d'un coup.
\end{remarque}

\section{Parcourir un tableau indexé}

\subsection{Avec la boucle \texttt{for}}

Comme les indices vont de \code{0} à \code{count(\$t) - 1}, la boucle \code{for}
est naturelle. La fonction \code{count} donne le \textbf{nombre d'éléments}.

\begin{codephp}
<?php
  $eleves = ["Awa", "Bello", "Chimene", "Daniel"];
  $n = count($eleves);          // n vaut 4

  for ($i = 0; $i < $n; $i++) {
    echo ($i + 1) . ". " . $eleves[$i] . "<br>";
  }
?>
\end{codephp}

\begin{exemple}
Le code ci-dessus affiche :\\
\code{1. Awa}\quad \code{2. Bello}\quad \code{3. Chimene}\quad \code{4. Daniel}
(chacun sur sa ligne grâce à \code{<br>}).
\end{exemple}

\subsection{Avec la boucle \texttt{foreach}}

\code{foreach} parcourt \textbf{toutes} les valeurs sans se soucier des indices :
c'est la façon la plus simple et la plus sûre de lire un tableau.

\begin{syntaxe}
\code{foreach (\$tableau as \$valeur) \{ \dots \}}
\end{syntaxe}

\begin{codephp}
<?php
  $villes = ["Yaoundé", "Douala", "Maroua"];

  foreach ($villes as $ville) {
    echo "Ville : " . $ville . "<br>";
  }
?>
\end{codephp}

\begin{remarque}
Avec \code{foreach}, pas besoin de gérer un compteur ni de connaître la taille du
tableau : on est sûr de ne sortir ni avant le début ni après la fin.
\end{remarque}

\section{Les tableaux associatifs}

Dans un tableau indexé, les clés sont des numéros. Dans un tableau
\textbf{associatif}, on remplace ces numéros par des \textbf{clés} de notre choix
(le plus souvent des mots), associées chacune à une valeur.

\begin{definition}
Un \textbf{tableau associatif} relie des \textbf{clés} (\emph{keys}) à des
\textbf{valeurs} (\emph{values}) grâce à la flèche \code{=>}. La clé sert
d'« étiquette » pour retrouver la valeur.
\end{definition}

\begin{codephp}
<?php
  // clé    =>  valeur
  $eleve = [
    "nom"     => "Awa Mballa",
    "classe"  => "Tle TI",
    "moyenne" => 14.5,
    "ville"   => "Garoua"
  ];

  echo $eleve["nom"];       // affiche : Awa Mballa
  echo $eleve["moyenne"];   // affiche : 14.5
?>
\end{codephp}

\begin{syntaxe}
\code{\$t = ["cle1" => valeur1, "cle2" => valeur2, \dots];}\\
Accès : \code{\$t["cle1"]}
\end{syntaxe}

\subsection{Parcourir un tableau associatif}

Pour obtenir à la fois la clé \textbf{et} la valeur, on utilise la forme complète
de \code{foreach} avec \code{=>}.

\begin{syntaxe}
\code{foreach (\$tableau as \$cle => \$valeur) \{ \dots \}}
\end{syntaxe}

\begin{codephp}
<?php
  $eleve = [
    "nom"     => "Awa Mballa",
    "classe"  => "Tle TI",
    "moyenne" => 14.5
  ];

  foreach ($eleve as $cle => $valeur) {
    echo $cle . " : " . $valeur . "<br>";
  }
  // Affiche :
  // nom : Awa Mballa
  // classe : Tle TI
  // moyenne : 14.5
?>
\end{codephp}

\begin{exemple}
On peut bâtir un \textbf{répertoire} reliant le nom d'un élève à sa moyenne, puis
le parcourir pour afficher chaque résultat :
\end{exemple}

\begin{codephp}
<?php
  $moyennes = [
    "Awa"     => 14.5,
    "Bello"   => 9.0,
    "Chimene" => 16.25,
    "Daniel"  => 11.75
  ];

  foreach ($moyennes as $nom => $moy) {
    echo $nom . " a obtenu " . $moy . "/20<br>";
  }
?>
\end{codephp}

\section{Les tableaux à plusieurs dimensions}

Une case d'un tableau peut elle-même contenir\dots\ un autre tableau\,! On obtient
un \textbf{tableau multidimensionnel}, très pratique pour représenter une
\textbf{liste de fiches} (une ligne par élève, plusieurs colonnes par ligne).

\begin{definition}
Un \textbf{tableau à deux dimensions} est un tableau dont chaque valeur est
elle-même un tableau. On accède à un élément avec \textbf{deux} indices ou clés :
\code{\$t[i][j]}.
\end{definition}

\begin{codephp}
<?php
  // Une fiche par élève (tableau de tableaux associatifs)
  $classe = [
    ["nom" => "Awa",     "moyenne" => 14.5],
    ["nom" => "Bello",   "moyenne" => 9.0],
    ["nom" => "Chimene", "moyenne" => 16.25]
  ];

  // Accès : 1re fiche, clé "nom"
  echo $classe[0]["nom"];       // affiche : Awa
  echo $classe[2]["moyenne"];   // affiche : 16.25

  // Parcours complet
  foreach ($classe as $fiche) {
    echo $fiche["nom"] . " : " . $fiche["moyenne"] . "/20<br>";
  }
?>
\end{codephp}

\begin{remarque}
Cette structure « tableau de fiches » est exactement ce que renverra une requête
de base de données MySQL (une ligne = une fiche). La maîtriser ici prépare le
chapitre sur PHP/MySQL.
\end{remarque}

\section{Fonctions utiles sur les tableaux}

PHP fournit de nombreuses fonctions toutes prêtes. Voici les plus courantes au
programme.

\begin{center}\small
\begin{tabular}{@{}lll@{}}
\toprule
\textbf{Fonction} & \textbf{Rôle} & \textbf{Exemple}\\
\midrule
\code{count(\$t)}        & nombre d'éléments      & \code{count([4,5,6])} → \code{3}\\
\code{sort(\$t)}         & trie en ordre croissant & \code{[3,1,2]} → \code{[1,2,3]}\\
\code{rsort(\$t)}        & trie en ordre décroissant & \code{[3,1,2]} → \code{[3,2,1]}\\
\code{in\_array(\$v,\$t)} & \$v est-il dans \$t ?   & renvoie \code{true}/\code{false}\\
\code{array\_push(\$t,\$v)} & ajoute \$v à la fin  & comme \code{\$t[] = \$v;}\\
\code{implode(\$sep,\$t)} & colle en une chaîne    & \code{implode(", ",\$t)}\\
\code{explode(\$sep,\$s)} & découpe une chaîne en tableau & \code{explode(",",\$s)}\\
\bottomrule
\end{tabular}
\end{center}

\begin{codephp}
<?php
  $notes = [12, 8, 15, 9, 17];

  echo count($notes);            // 5

  sort($notes);                  // $notes devient [8, 9, 12, 15, 17]
  rsort($notes);                 // $notes devient [17, 15, 12, 9, 8]

  if (in_array(15, $notes)) {
    echo "La note 15 existe.";
  }

  array_push($notes, 20);        // ajoute 20 à la fin

  // implode : tableau -> chaîne
  $eleves = ["Awa", "Bello", "Chimene"];
  echo implode(" - ", $eleves);  // affiche : Awa - Bello - Chimene

  // explode : chaîne -> tableau
  $liste = "Yaoundé;Douala;Maroua";
  $villes = explode(";", $liste); // ["Yaoundé", "Douala", "Maroua"]
  echo $villes[1];               // affiche : Douala
?>
\end{codephp}

\begin{attention}
\code{sort} et \code{rsort} \textbf{modifient directement} le tableau passé en
argument (on ne récupère pas le résultat avec un \code{=}). Après
\code{sort(\$notes);}, c'est \code{\$notes} lui-même qui est trié.
\end{attention}

\section{Les fonctions}

Quand un même calcul revient souvent (calculer une moyenne, décider d'une
admission\dots), on l'enferme dans une \textbf{fonction} que l'on appelle ensuite
autant de fois que nécessaire. C'est la clé d'un code clair et réutilisable.

\begin{definition}
Une \textbf{fonction} est un bloc d'instructions portant un \textbf{nom}. On la
\textbf{déclare} une fois, puis on l'\textbf{appelle} par son nom. Elle peut
recevoir des \textbf{paramètres} (données d'entrée) et \textbf{renvoyer} un
résultat avec \code{return}.
\end{definition}

\subsection{Déclarer et appeler une fonction}

\begin{syntaxe}
\code{function nom(\$param1, \$param2) \{ \dots\ return \$resultat; \}}
\end{syntaxe}

\begin{codephp}
<?php
  // Déclaration : on définit la fonction
  function carre($x) {
    return $x * $x;
  }

  // Appel : on l'utilise
  echo carre(5);        // affiche : 25
  $r = carre(9);        // $r vaut 81
  echo $r;
?>
\end{codephp}

\begin{remarque}
\code{return} \textbf{renvoie} une valeur \emph{et arrête} la fonction : les
instructions placées après ne sont pas exécutées. Une fonction sans \code{return}
peut tout de même travailler (par exemple afficher quelque chose avec
\code{echo}).
\end{remarque}

\subsection{Paramètres par défaut}

On peut donner une valeur \textbf{par défaut} à un paramètre : si l'appel ne la
précise pas, cette valeur est utilisée.

\begin{codephp}
<?php
  // La devise vaut "FCFA" si on ne la précise pas
  function prix($montant, $devise = "FCFA") {
    return $montant . " " . $devise;
  }

  echo prix(2500);          // affiche : 2500 FCFA
  echo prix(10, "euros");   // affiche : 10 euros
?>
\end{codephp}

\subsection{Portée des variables}

\begin{definition}
La \textbf{portée} (\emph{scope}) d'une variable est la zone du programme où elle
existe. Une variable créée \textbf{dans} une fonction est \textbf{locale} : elle
n'existe pas en dehors. Une variable créée en dehors est \textbf{globale} et
n'est, par défaut, \textbf{pas visible} à l'intérieur d'une fonction.
\end{definition}

\begin{codephp}
<?php
  $tva = 0.1925;   // variable globale

  function montantTTC($ht) {
    global $tva;            // on rend $tva visible ici
    return $ht * (1 + $tva);
  }

  echo montantTTC(1000);    // affiche : 1192.5
?>
\end{codephp}

\begin{attention}
Sans le mot-clé \code{global}, la variable \code{\$tva} serait \code{NULL} à
l'intérieur de la fonction. Mieux encore : on évite \code{global} en
\textbf{passant} la donnée en paramètre, ce qui rend la fonction plus claire.
\end{attention}

\section{Quelques fonctions PHP utiles}

\subsection{Sur les chaînes de caractères}

\begin{center}\small
\begin{tabular}{@{}lll@{}}
\toprule
\textbf{Fonction} & \textbf{Rôle} & \textbf{Exemple}\\
\midrule
\code{strlen(\$s)}            & longueur (nb de caractères) & \code{strlen("Awa")} → \code{3}\\
\code{strtoupper(\$s)}        & tout en majuscules     & → \code{"AWA"}\\
\code{strtolower(\$s)}        & tout en minuscules     & → \code{"awa"}\\
\code{substr(\$s,\$d,\$l)}    & extrait une portion    & \code{substr("Yaoundé",0,3)} → \code{"Yao"}\\
\code{str\_replace(\$a,\$b,\$s)} & remplace \$a par \$b & remplace dans \$s\\
\bottomrule
\end{tabular}
\end{center}

\begin{codephp}
<?php
  $nom = "Awa Mballa";

  echo strlen($nom);                 // 10 (l'espace compte)
  echo strtoupper($nom);             // AWA MBALLA
  echo strtolower($nom);             // awa mballa
  echo substr($nom, 0, 3);           // Awa (3 caractères dès l'indice 0)
  echo str_replace("a", "@", $nom);  // Aw@ Mb@ll@
?>
\end{codephp}

\subsection{Sur les nombres}

\begin{center}\small
\begin{tabular}{@{}lll@{}}
\toprule
\textbf{Fonction} & \textbf{Rôle} & \textbf{Exemple}\\
\midrule
\code{round(\$x,\$d)}  & arrondit à \$d décimales & \code{round(14.567,2)} → \code{14.57}\\
\code{abs(\$x)}        & valeur absolue           & \code{abs(-5)} → \code{5}\\
\code{max(\$a,\$b,\dots)} & le plus grand         & \code{max(8,15,9)} → \code{15}\\
\code{min(\$a,\$b,\dots)} & le plus petit         & \code{min(8,15,9)} → \code{8}\\
\code{rand(\$a,\$b)}   & entier aléatoire \$a..\$b & \code{rand(1,6)} (dé)\\
\bottomrule
\end{tabular}
\end{center}

\begin{codephp}
<?php
  echo round(14.567, 2);   // 14.57
  echo abs(-8);            // 8
  echo max(8, 15, 9);     // 15
  echo min(8, 15, 9);     // 8
  echo rand(1, 6);        // un entier au hasard entre 1 et 6

  // max et min acceptent aussi un tableau
  $notes = [12, 8, 17, 9];
  echo max($notes);       // 17
?>
\end{codephp}

\section{Exemples complets}

\begin{exemple}
\textbf{Moyenne d'un tableau de notes.} On parcourt le tableau pour cumuler la
somme, puis on divise par le nombre de notes.
\end{exemple}

\begin{codephp}
<?php
  $notes = [12, 8, 15, 9, 17];
  $somme = 0;

  foreach ($notes as $note) {
    $somme = $somme + $note;
  }

  $moyenne = $somme / count($notes);
  echo "Moyenne de la classe : " . round($moyenne, 2) . "/20";
  // affiche : Moyenne de la classe : 12.2/20
?>
\end{codephp}

\begin{exemple}
\textbf{Fonction \code{estAdmis}.} Elle reçoit une moyenne et renvoie \code{true}
si l'élève est admis (moyenne $\geq 10$), \code{false} sinon.
\end{exemple}

\begin{codephp}
<?php
  function estAdmis($moyenne) {
    return $moyenne >= 10;
  }

  $moy = 11.5;
  if (estAdmis($moy)) {
    echo "Admis";
  } else {
    echo "Non admis";
  }
?>
\end{codephp}

\begin{exemple}
\textbf{Fonction qui calcule une moyenne.} On regroupe le calcul précédent dans
une fonction réutilisable, qui reçoit un tableau et renvoie la moyenne.
\end{exemple}

\begin{codephp}
<?php
  function moyenne($notes) {
    $somme = 0;
    foreach ($notes as $n) {
      $somme = $somme + $n;
    }
    return $somme / count($notes);
  }

  $classeA = [12, 8, 15, 9, 17];
  $classeB = [10, 11, 14];

  echo round(moyenne($classeA), 2);  // 12.2
  echo round(moyenne($classeB), 2);  // 11.67
?>
\end{codephp}

\begin{remarque}
Une fois \code{moyenne} écrite, on la réutilise pour n'importe quelle classe sans
réécrire la boucle. C'est tout l'intérêt des fonctions : \textbf{écrire une fois,
réutiliser partout}.
\end{remarque}

\section{Exercices}

\rubrique{Application}

\exo{Afficher un tableau}\diff{1}
On donne \code{\$villes = ["Yaoundé", "Douala", "Maroua", "Bafoussam"]}. Affiche
chaque ville précédée de son numéro de position (1, 2, 3, 4), une par ligne.

\exo{Max, min, somme et moyenne}\diff{1}
Soit \code{\$notes = [12, 8, 15, 9, 17, 6]}. Affiche la plus grande note, la plus
petite, la somme totale et la moyenne arrondie à deux décimales.

\exo{Rechercher un élève}\diff{1}
On dispose de \code{\$classe = ["Awa", "Bello", "Chimene", "Daniel"]}. Écris un
script qui indique si l'élève « Chimene » figure dans la liste, puis si « Eric »
y figure.

\exo{Compter les admis}\diff{2}
Soit le tableau associatif suivant :
\begin{codephp}
$moyennes = ["Awa"=>14, "Bello"=>8, "Chimene"=>16, "Daniel"=>9, "Eric"=>11];
\end{codephp}
Compte combien d'élèves sont admis (moyenne $\geq 10$) et affiche le résultat.

\exo{Fonction \texttt{mention}}\diff{2}
Écris une fonction \code{mention(\$m)} qui renvoie : « Excellent » si
$m \geq 16$, « Bien » si $14 \leq m < 16$, « Assez bien » si $12 \leq m < 14$,
« Passable » si $10 \leq m < 12$ et « Insuffisant » sinon. Teste-la sur quelques
moyennes.

\rubrique{Entraînement}

\exo{Inverser un tableau}\diff{2}
Sans utiliser la fonction toute prête \code{array\_reverse}, écris un script qui
construit un nouveau tableau contenant les éléments de \code{\$t = [10, 20, 30,
40]} dans l'ordre inverse, puis l'affiche.

\exo{Bulletin de la classe}\diff{2}
Soit le tableau de fiches :
\begin{codephp}
$classe = [["nom"=>"Awa","moy"=>14], ["nom"=>"Bello","moy"=>8], ["nom"=>"Chimene","moy"=>16]];
\end{codephp}
Affiche pour chaque élève son nom, sa moyenne et sa mention (réutilise la
fonction \code{mention} de l'exercice 5).

\exo{Fréquences des notes}\diff{3}
On donne \code{\$notes = [12, 8, 12, 15, 8, 12, 9]}. Construis un tableau
associatif qui compte combien de fois chaque note apparaît, puis affiche le
résultat sous la forme \code{12 -> 3 fois}.

\exo{Initiales}\diff{3}
Écris une fonction \code{initiales(\$nomComplet)} qui, à partir d'une chaîne
de noms séparés par des espaces, renvoie les initiales en majuscules. Par
exemple, \code{"Awa Mballa Ngono"} doit donner \code{"A.M.N."}.
Indice : utilise \code{explode}, puis \code{substr} et \code{strtoupper}.

\section{Corrigés}

\corrige{de l'exercice 1}
\begin{codephp}
<?php
  $villes = ["Yaoundé", "Douala", "Maroua", "Bafoussam"];

  for ($i = 0; $i < count($villes); $i++) {
    echo ($i + 1) . ". " . $villes[$i] . "<br>";
  }
  // On pouvait aussi utiliser foreach avec un compteur séparé.
?>
\end{codephp}

\corrige{de l'exercice 2}
\begin{codephp}
<?php
  $notes = [12, 8, 15, 9, 17, 6];

  $somme = 0;
  foreach ($notes as $n) {
    $somme = $somme + $n;
  }

  echo "Maximum : " . max($notes) . "<br>";
  echo "Minimum : " . min($notes) . "<br>";
  echo "Somme : " . $somme . "<br>";
  echo "Moyenne : " . round($somme / count($notes), 2);
?>
\end{codephp}

\corrige{de l'exercice 3}
\begin{codephp}
<?php
  $classe = ["Awa", "Bello", "Chimene", "Daniel"];

  if (in_array("Chimene", $classe)) {
    echo "Chimene est dans la classe.<br>";
  } else {
    echo "Chimene est absente.<br>";
  }

  if (in_array("Eric", $classe)) {
    echo "Eric est dans la classe.";
  } else {
    echo "Eric n'est pas dans la classe.";
  }
?>
\end{codephp}

\corrige{de l'exercice 4}
\begin{codephp}
<?php
  $moyennes = [
    "Awa"     => 14,
    "Bello"   => 8,
    "Chimene" => 16,
    "Daniel"  => 9,
    "Eric"    => 11
  ];

  $nbAdmis = 0;
  foreach ($moyennes as $nom => $moy) {
    if ($moy >= 10) {
      $nbAdmis = $nbAdmis + 1;
    }
  }

  echo "Nombre d'admis : " . $nbAdmis . " sur " . count($moyennes);
  // affiche : Nombre d'admis : 3 sur 5
?>
\end{codephp}

\corrige{de l'exercice 5}
\begin{codephp}
<?php
  function mention($m) {
    if ($m >= 16) {
      return "Excellent";
    } elseif ($m >= 14) {
      return "Bien";
    } elseif ($m >= 12) {
      return "Assez bien";
    } elseif ($m >= 10) {
      return "Passable";
    } else {
      return "Insuffisant";
    }
  }

  echo mention(17);   // Excellent
  echo mention(13);   // Assez bien
  echo mention(7);    // Insuffisant
?>
\end{codephp}

\begin{remarque}
On teste les seuils du plus grand au plus petit. Dès qu'une condition est vraie,
\code{return} renvoie la mention et arrête la fonction : inutile de re-tester
$m < 16$, c'est déjà garanti par l'ordre des \code{elseif}.
\end{remarque}

\corrige{de l'exercice 6}
\begin{codephp}
<?php
  $t = [10, 20, 30, 40];
  $inverse = [];

  for ($i = count($t) - 1; $i >= 0; $i--) {
    $inverse[] = $t[$i];   // on ajoute du dernier vers le premier
  }

  // Affichage
  foreach ($inverse as $valeur) {
    echo $valeur . " ";    // affiche : 40 30 20 10
  }
?>
\end{codephp}

\corrige{de l'exercice 7}
\begin{codephp}
<?php
  function mention($m) {
    if ($m >= 16)      return "Excellent";
    elseif ($m >= 14)  return "Bien";
    elseif ($m >= 12)  return "Assez bien";
    elseif ($m >= 10)  return "Passable";
    else               return "Insuffisant";
  }

  $classe = [
    ["nom" => "Awa",     "moy" => 14],
    ["nom" => "Bello",   "moy" => 8],
    ["nom" => "Chimene", "moy" => 16]
  ];

  foreach ($classe as $eleve) {
    echo $eleve["nom"] . " : " . $eleve["moy"]
       . "/20 (" . mention($eleve["moy"]) . ")<br>";
  }
  // Awa : 14/20 (Bien)
  // Bello : 8/20 (Insuffisant)
  // Chimene : 16/20 (Excellent)
?>
\end{codephp}

\corrige{de l'exercice 8}
\begin{codephp}
<?php
  $notes = [12, 8, 12, 15, 8, 12, 9];
  $freq = [];   // tableau associatif note => nombre

  foreach ($notes as $n) {
    if (in_array($n, array_keys($freq))) {
      $freq[$n] = $freq[$n] + 1;   // déjà vue : on incrémente
    } else {
      $freq[$n] = 1;               // première fois
    }
  }

  foreach ($freq as $note => $nombre) {
    echo $note . " -> " . $nombre . " fois<br>";
  }
  // 12 -> 3 fois
  // 8 -> 2 fois
  // 15 -> 1 fois
  // 9 -> 1 fois
?>
\end{codephp}

\corrige{de l'exercice 9}
\begin{codephp}
<?php
  function initiales($nomComplet) {
    $mots = explode(" ", $nomComplet);   // découpe sur les espaces
    $resultat = "";

    foreach ($mots as $mot) {
      $premiere = substr($mot, 0, 1);    // 1re lettre du mot
      $resultat = $resultat . strtoupper($premiere) . ".";
    }

    return $resultat;
  }

  echo initiales("Awa Mballa Ngono");   // affiche : A.M.N.
?>
\end{codephp}

\begin{aretenir}
Un \textbf{tableau} regroupe plusieurs valeurs : indices numériques (départ à
\code{0}) ou clés \code{=>} valeurs. On le parcourt avec \code{foreach}
(\code{\$t as \$cle => \$val}) et on l'exploite avec \code{count}, \code{sort},
\code{in\_array}\dots\ Une \textbf{fonction} (\code{function nom(\$p) \{ return
\dots; \}}) enferme un calcul pour l'\textbf{écrire une fois et le réutiliser
partout}.
\end{aretenir}
